\documentclass{article}
\usepackage[T1]{fontenc}
\usepackage{lmodern}
\usepackage{graphicx}
\usepackage{amsmath,amssymb}
\usepackage[a4paper,margin=2cm]{geometry}
\usepackage[absolute,overlay]{textpos}
\setlength{\TPHorizModule}{1mm}
\setlength{\TPVertModule}{1mm}
\usepackage[indonesian]{babel}
\usepackage{hyperref}
\setlength{\emergencystretch}{2em}
% unit: Halaman 9

\begin{document}

% === Halaman 9 (llm) ===

\begin{itemize}
    \item Steganalisis diperlukan di dalam \textit{forensic image analysis}
    \item \textbf{Forensic Image Analysis} is the application of image science and domain expertise to interpret the \underline{content of} an image and/or the image itself in legal matters.
    \item Subdisiplin dari \textit{Forensic Image Analysis}:
    \begin{enumerate}
        \item Photogrammetry
        \item Photographic Comparison
        \item \underline{Content Analysis}
        \item \underline{Image Authentication}
    \end{enumerate}
\end{itemize}

\end{document}
